% BHU PYQ -- Transcribed & Corrected 100% Accurate
% Paper: PHYMJ21 / PHYMN21: Thermal Physics
% Exam: B.Sc. (Hons.) Semester II Examination 2024-25 (NEP Mid-Term)
% Department: Department of Physics

\documentclass[11pt,a4paper]{article}
\usepackage[a4paper,top=2.5cm,bottom=2.5cm,left=2.5cm,right=2.5cm,headheight=14pt]{geometry}
\usepackage{amsmath,amssymb}
\usepackage[T1]{fontenc}
\usepackage[utf8]{inputenc}
\usepackage{lmodern,microtype,fancyhdr,enumitem,hyperref,lastpage,parskip}

\hypersetup{colorlinks=true,urlcolor=black,linkcolor=black,pdftitle={PHYMJ21 / PHYMN21: Thermal Physics -- BHU Sem II Mid-Term 2024-25 (NEP)}}
\pagestyle{fancy}\fancyhf{}
\fancyhead[L]{\small\textit{Banaras Hindu University}}
\fancyhead[R]{\small\textit{PHYMJ21 / PHYMN21: Thermal Physics (Mid-Term)}}
\fancyfoot[C]{\small Page \thepage\ of \pageref{LastPage}}

\newlist{parts}{enumerate}{2}
\setlist[parts,1]{label=(\roman*),leftmargin=*,itemsep=4pt,topsep=2pt}
\newcommand{\pts}[1]{\hfill\textit{[#1]}}

\begin{document}

\begin{center}
  {\large\bfseries BANARAS HINDU UNIVERSITY}\\[2pt]
  {\normalsize B.Sc. (Hons.) Semester II Examination 2024-25 (NEP)}\\[6pt]
  \rule{0.8\linewidth}{0.5pt}\\[6pt]
  {\large\bfseries DEPARTMENT OF PHYSICS}\\[4pt]
  {\large\bfseries Mid-Semester Examination}\\[4pt]
  {\large\bfseries Paper Code: PHYMJ21 / PHYMN21 --- Thermal Physics}\\[4pt]
  {\normalsize Time Allowed: 1:00 Hour\hfill Maximum Marks: 20}
\end{center}

\rule{\linewidth}{0.4pt}\smallskip
\noindent\textit{Instruction: Attempt all questions.}
\bigskip

\noindent\textbf{Question 1}\pts{2} \\
Calculate the work done by $2\text{ moles}$ of an ideal gas in terms of gas constant $R$ whose volume expands isothermally to $4$ times of its initial volume.

\medskip\rule{\linewidth}{0.2pt}\medskip

\noindent\textbf{Question 2} \textit{Write short notes on any two of the following:}\pts{$2.5 + 2.5 = 5$}
\begin{parts}
  \item Thermodynamic equilibrium
  \item Second law of thermodynamics
  \item Refrigerator
\end{parts}

\medskip\rule{\linewidth}{0.2pt}\medskip

\noindent\textbf{Question 3}\pts{3} \\
Two Carnot engines $X$ and $Y$ are operated in series. Engine $X$ receives heat from a source at temperature $900\text{ K}$ and rejects heat to a sink at temperature $T\text{ K}$. The second engine $Y$ receives heat at temperature $T\text{ K}$ and rejects heat to its sink at temperature $400\text{ K}$. For what value of temperature $T$ will the efficiencies of both engines be the same?

\medskip\rule{\linewidth}{0.2pt}\medskip

\noindent\textbf{Question 4}\pts{5} \\
Write Maxwell's law of speed distribution and by using this derive the relation for the most probable speed of gas molecules.

\medskip\rule{\linewidth}{0.2pt}\medskip

\noindent\textbf{Question 5}\pts{2} \\
Write the relationship between mean free path and gas temperature. Explain the terms used.

\medskip\rule{\linewidth}{0.2pt}\medskip

\noindent\textbf{Question 6}\pts{3} \\
If the diameter of a gas molecule is $2 \times 10^{-8}\text{ cm}$, find the value of mean free path at STP ($T = 273\text{ K}$, $P = 1.013 \times 10^5\text{ N/m}^2$, $k = 1.38 \times 10^{-23}\text{ J/K}$).

\vfill\rule{\linewidth}{0.4pt}
\begin{center}\small
\href{https://bhu-exam.vercel.app/}{Click here to visit our website}\\[4pt]
\href{https://me-aryan.vercel.app/}{Click here to visit our contributor}
\end{center}

\end{document}
